\uuid{2Fck}
\titre{Fonction croissante/décroissante}
\auteur{bodin ; hanani ; mzari}
\organisation{exo7}
\licence{CC BY-NC-SA 4.0}
\datecreate{2020-01-20}
\niveau{L1}
\module{Analyse}
\chapitre{Continuité, limite et étude de fonctions réelles}
\sousChapitre{Autre}
\difficulte{2}
\theme{Notion de fonctions, fonction croissante, décroissante}
\exo7qcmid{lille-1.continuite.006}

\contenu{
\begin{question}
\qclassification{123.00}
\qtags{motcle=fonction croissante/décroissante}
Quelles sont les assertions vraies ?
\begin{answers}
    \bad{La fonction $x \mapsto \frac{1}{x}$ est décroissante car sa dérivée $x \mapsto -\frac{1}{x^2}$ est partout négative.}

    \good{Une fonction périodique et croissante est constante.}

    \good{Si $f : \Rr \to ]0,+\infty[$ est croissante, alors $1/f$ est décroissante.}

    \bad{Si $f$ et $g$ sont croissantes, alors $f-g$ est croissante.}
\end{answers}
\begin{explanations}
La fonction $x \mapsto \frac{1}{x}$ est décroissante sur $]0,+\infty[$ et $]-\infty,0[$ mais pas sur $\Rr^*$.
\end{explanations}
\end{question}
}
